\chapter{La terminal sin miedo}\label{cap:18-terminal}

La \termino{terminal} es una ventana donde se le dan órdenes al equipo
escribiendo, en lugar de hacer clic. Para casi todo lo de este libro no
hace falta, pero a veces es el camino más corto, y es como suelen
responder quienes te ayudan.

Se abre desde \menu{Aplicaciones > Sistema > Terminal} (o el icono del
panel de abajo en XFCE).

\captura{terminal}{La terminal de XFCE}

\section{Cómo se usa}

\begin{itemize}
  \item Escribe la orden y pulsa \keys{Enter}.
  \item Las mayúsculas, las minúsculas y los espacios importan.
  \item \keys{\arrowkeyup} trae la orden anterior; así no hay que escribirla
        de nuevo.
  \item \keys{Tab} completa nombres de archivos y órdenes: escribe las
        primeras letras y pulsa \keys{Tab}.
  \item Para pegar, usa \keys{\ctrl + Mayús + V} o clic derecho,
        \menu{Pegar} (\keys{\ctrl + V} solo no funciona aquí).
  \item \keys{\ctrl + C} detiene la orden que está corriendo.
  \item Cuando una orden pide la contraseña, lo que escribes no se ve.
        Escríbela igual y pulsa \keys{Enter}.
\end{itemize}

\section{Órdenes útiles}

\begin{center}
\begin{tabularx}{\textwidth}{>{\ttfamily}lX}
\toprule
\normalfont Orden & Qué hace \\
\midrule
ls & Muestra los archivos de la carpeta actual \\
cd Documentos & Entra en la carpeta \emph{Documentos} \\
cd & Vuelve a la carpeta personal \\
pwd & Dice en qué carpeta estás \\
fastfetch & Resumen del equipo y del sistema \\
btop & Qué está usando el procesador y la memoria (\keys{Q} para salir) \\
df -h & Cuánto espacio queda en los discos \\
free -h & Cuánta memoria está en uso \\
ip addr & Las conexiones de red y sus direcciones \\
nmcli device wifi list & Las redes wifi que se ven \\
doas apk upgrade & Actualiza todo lo de Alpine \\
lazaro-diagnostico & El informe para pedir ayuda \\
\bottomrule
\end{tabularx}
\end{center}

\section{Ser administrador: doas}

Las órdenes que cambian el sistema se escriben con \texttt{doas} delante:
\begin{terminal}
doas apk add gimp
\end{terminal}
\texttt{doas} pide tu contraseña y ejecuta la orden con permisos de
administrador. Es el equivalente a \enquote{Ejecutar como administrador} de
Windows o al \texttt{sudo} de otras distribuciones. Úsalo solo con órdenes
que entiendas o que vengan de una fuente de confianza.

\begin{nota}
Para aprender más, la wiki de Alpine explica cada parte del sistema
\cite{alpine-wiki}, y el gestor de paquetes \texttt{apk} tiene su propia
página \cite{alpine-apk}.
\end{nota}

\fuentes
